\chapter{Architettura generale}
\label{ch:architettura}

\section{Visione}

L'applicazione è organizzata in tre aree funzionali indipendenti
(\module{Editor}, \module{Scheduler}, \module{Simulator}) che
comunicano esclusivamente tramite tipi in \repo{Domain/Model/}.
Le dipendenze puntano sempre verso il dominio (regola delle frecce).

\section{Struttura repository}

\begin{figure}[htbp]
\centering
\begin{tikzpicture}[
  box/.style={draw, rounded corners, minimum width=3.2cm,
    minimum height=0.9cm, align=center, font=\small},
  arr/.style={-{Stealth[length=2mm]}, thick}
]
  \node[box, fill=blue!8] (ui) {UI / ViewModel};
  \node[box, fill=green!10, below=1.2cm of ui] (svc) {Services};
  \node[box, fill=orange!12, below=1.2cm of svc] (dom) {Domain/Model};
  \node[box, fill=gray!15, right=2.8cm of svc] (infra) {Infrastructure};
  \node[box, fill=gray!15, left=2.8cm of svc] (models) {Models (bridge)};

  \draw[arr] (ui) -- (svc);
  \draw[arr] (svc) -- (dom);
  \draw[arr] (infra) -- (dom);
  \draw[arr] (models) -- (dom);
  \draw[arr, dashed] (ui) -- (models);
\end{tikzpicture}
\caption{Direzione delle dipendenze (semplificata).}
\end{figure}

\section{Refactor Strangler (Fasi 0--4)}

Il monolite \texttt{RailwayScheduleOptimizer} è stato sostituito da
\repo{Services/Scheduling/}. La migrazione dei modelli da
\repo{Models/} a \repo{Domain/Model/} è in corso: i tipi infrastrutturali
(\texttt{Node}, \texttt{Edge}, \texttt{Signal}, \ldots) risiedono già
nel dominio; \texttt{NetworkModel} e i DTO restano nel layer bridge.

\section{Swift Packages (scaffold)}

\begin{table}[htbp]
\centering
\caption{Pacchetti SPM locali (\repo{Packages/}).}
\begin{tabular}{lll}
  \toprule
  \textbf{Package} & \textbf{Stato} & \textbf{Note} \\
  \midrule
  \texttt{FDCDomain}   & Compilabile & 9 tipi Foundation via symlink \\
  \texttt{FDCScheduling} & Scaffold & Non linkato al target Xcode \\
  \bottomrule
\end{tabular}
\end{table}

\section{Test}

\begin{itemize}
  \item \textbf{20 test XCTest} nel target \texttt{FdC Railway ManagerTests}.
  \item \textbf{2 test} nel package \texttt{FDCDomain}.
  \item I servizi di scheduling sono \texttt{struct} pure su
    \texttt{RailwayTopology} (no \texttt{@MainActor}, no singleton).
\end{itemize}

\section{Feature flags}

Funzionalità sperimentali in \repo{FeatureFlags.swift}:
\texttt{simulatorGameMode}, \texttt{realtimeConflictHighlight}.
Il flag \texttt{useLegacyScheduleOptimizer} è stato rimosso con la Fase 4.
